\documentclass[11pt,a4paper]{article}

\usepackage[utf8]{inputenc}
\usepackage[margin=0.85in]{geometry}
\usepackage{amsmath,amssymb}
\usepackage{enumitem}
\usepackage{booktabs}
\usepackage{tcolorbox}
\usepackage{hyperref}
\usepackage{xcolor}
\usepackage{fancyhdr}

\hypersetup{
    colorlinks=true,
    linkcolor=blue!75!black,
    urlcolor=blue!75!black
}

% Custom boxes for Q&A
\newtcolorbox{answerbox}{
    colback=blue!5!white,
    colframe=blue!60!black,
    title=\textbf{Ideal Answer (Plain Language)},
    arc=3mm,
    fonttitle=\bfseries
}

\newtcolorbox{conceptbox}[1]{
    colback=orange!5!white,
    colframe=orange!70!black,
    title=\textbf{#1},
    arc=3mm,
    fonttitle=\bfseries
}

\newtcolorbox{tipbox}{
    colback=green!5!white,
    colframe=green!60!black,
    title=\textbf{Viva Tip for the Team},
    arc=3mm,
    fonttitle=\bfseries
}

\pagestyle{fancy}
\fancyhf{}
\rhead{GAI-PLA -- Viva Preparation Notes}
\lhead{University Capstone Project}
\cfoot{\thepage}

\begin{document}

% -----------------------------------------------------------------------------
% Title Header
% -----------------------------------------------------------------------------
\begin{center}
    {\huge \textbf{GAI-PLA: Viva Preparation Guide}}\\[0.2cm]
    {\large \textbf{Generative AI-Based Personalized Learning Agent}}\\[0.3cm]
    \textit{Plain-Language Explanations, Core Concepts, and Evaluator Q\&A Cheat Sheet}\\[0.2cm]
    \textbf{Team Author:} Anurag Gaur \quad $\vert$ \quad \textbf{Repository:} \href{https://github.com/anuraggaur29/GAI-PLA}{github.com/anuraggaur29/GAI-PLA}
\end{center}
\vspace{0.3cm}
\hrule
\vspace{0.4cm}

% -----------------------------------------------------------------------------
% Section 1: 30-Second Elevator Pitch
% -----------------------------------------------------------------------------
\section{The 30-Second Elevator Pitch (Master This First)}

\begin{conceptbox}{What is GAI-PLA in simple terms?}
``GAI-PLA is an intelligent AI tutoring system. Unlike ChatGPT, which just gives long text explanations when you ask a question, GAI-PLA acts like a real private teacher:
\begin{enumerate}[noitemsep]
    \item You give it any topic (e.g., \textit{Pythagoras Theorem}).
    \item It breaks down the topic into a visual \textbf{Knowledge Graph} of sub-concepts.
    \item It gives you an \textbf{Adaptive Quiz} to test what you know and what you don't.
    \item It diagnoses your exact weak points (e.g., \textit{``You understand the theorem, but struggle with Right Triangle basics''}).
    \item Finally, it generates a \textbf{custom lesson} and downloadable \textbf{PowerPoint slides} focusing \textit{only} on what you got wrong, using personalized teaching styles like the Feynman Technique.''
\end{enumerate}
\end{conceptbox}

% -----------------------------------------------------------------------------
% Section 2: The Big "Why" Questions (Evaluator Favorites)
% -----------------------------------------------------------------------------
\section{The Big ``Why'' Questions (Evaluator Favorites)}

\subsection*{Q1. Why not just use ChatGPT or Claude directly? Why do we need GAI-PLA?}
\begin{answerbox}
\textbf{The problem with raw LLMs:}
\begin{itemize}[noitemsep]
    \item \textbf{No diagnostic assessment:} ChatGPT just dumps an essay; it doesn't test if the student actually understands.
    \item \textbf{No structured learning loop:} A student doesn't know where to start or what prerequisite they are missing.
    \item \textbf{One-size-fits-all:} ChatGPT explains everything equally instead of focusing 80\% of the time on the 20\% the student struggled with.
\end{itemize}
\textbf{How GAI-PLA solves this:}
GAI-PLA wraps the AI into a strict pedagogical workflow: \textbf{Assess first $\to$ Diagnose weaknesses $\to$ Generate targeted remediation $\to$ Export slides}. It turns an AI from a ``talking encyclopedia'' into a ``structured personal tutor.''
\end{answerbox}

\subsection*{Q2. What is LangGraph, and why did you use it instead of a simple Python script?}
\begin{answerbox}
\begin{itemize}[noitemsep]
    \item \textbf{Simple Analogy:} If regular code is a straight conveyor belt, LangGraph is a smart train network with switches and tracks.
    \item \textbf{Why we use it:} Our application has multiple distinct steps (Chatting, Knowledge Graph, Quizzing, Grading, Lessons, Slides). LangGraph maintains the \textbf{Student State} (what questions were answered, which nodes were failed) across all steps.
    \item If the user sends a message, an \textbf{Intent Node} decides whether to chat, create a quiz, or build slides, keeping the application reliable and modular without messy \texttt{if-else} spaghetti code.
\end{itemize}
\end{answerbox}

\subsection*{Q3. What is the role of ``Instructor'' and Pydantic?}
\begin{answerbox}
\begin{itemize}[noitemsep]
    \item \textbf{The Problem:} LLMs output free-form text. If an LLM returns bad JSON with a missing bracket, the frontend website crashes.
    \item \textbf{The Solution (Instructor + Pydantic):} Instructor forces the LLM to strictly follow defined blueprints (Pydantic models). It guarantees that when we ask for a quiz question, the LLM \textit{always} returns exactly 4 options, a question text, and 1 correct answer. If the LLM makes a formatting mistake, Instructor automatically retries and fixes it before sending it to the user.
\end{itemize}
\end{answerbox}

% -----------------------------------------------------------------------------
% Section 3: Step-by-Step Architecture (The 6 Steps)
% -----------------------------------------------------------------------------
\section{How the System Works Step-by-Step (The Pipeline)}

When the professor asks: \textit{``Walk me through what happens when a user types a topic,''} explain these 6 steps:

\begin{enumerate}[leftmargin=*]
    \item \textbf{Step 1: Intent Routing (\texttt{src/agent/nodes/intent.py})}
    The user enters \textit{``Pythagoras theorem''}. The intent node analyzes the prompt and switches the system into \textbf{Learning Mode}.
    
    \item \textbf{Step 2: Knowledge Graph Synthesis (\texttt{knowledge\_graph.py})}
    The AI breaks the topic into 7 foundational nodes (e.g., \textit{Right Triangle, Squares and Square Roots, Pythagorean Theorem, Distance Formula}). These nodes are rendered interactively on the frontend using Vis.js.
    
    \item \textbf{Step 3: Adaptive Diagnostic Quiz (\texttt{quiz.py})}
    GAI-PLA generates 14 calibrated multiple-choice questions mapped directly to those 7 knowledge points (2 questions per node) to establish the learner's baseline.
    
    \item \textbf{Step 4: Grading \& Weakness Diagnosis (\texttt{grading.py} \& \texttt{diagnosis.py})}
    The student submits answers. The system grades each response, calculates overall accuracy (e.g., 50\%), and tags each node as either \textbf{Mastered} or \textbf{Needs Review}.
    
    \item \textbf{Step 5: Targeted Remediation (\texttt{lesson.py} \& \texttt{slides.py})}
    Instead of re-explaining the entire subject, GAI-PLA generates:
    \begin{itemize}[noitemsep]
        \item A personalized \textbf{Markdown lesson} streamed in real-time using Server-Sent Events (SSE).
        \item A 16:9 \textbf{Presentation Slide Deck} that can be previewed in the browser or downloaded as an editable \texttt{.pptx} PowerPoint file via \texttt{python-pptx}.
    \end{itemize}
    
    \item \textbf{Step 6: L2 Reflection \& Memory (\texttt{review.py} + \texttt{student\_memory/})}
    The system summarizes mistakes into structured reflection notes stored in the vector database with an Ebbinghaus forgetting curve, so next week the AI remembers what the student struggled with.
\end{enumerate}

% -----------------------------------------------------------------------------
% Section 4: Key Modules & Technologies Cheat Sheet
% -----------------------------------------------------------------------------
\section{Tech Stack Quick Reference (Cheat Sheet)}

\begin{table}[htbp]
\centering
\small
\begin{tabular}{lll}
\toprule
\textbf{Component} & \textbf{Technology Used} & \textbf{Why it was chosen} \\
\midrule
Backend API & FastAPI + Uvicorn & Async high-performance REST endpoints + SSE streaming. \\
Workflow Engine & LangGraph & Manages state transitions and pedagogical routing. \\
LLM Model & Google Gemini Flash & High speed, low latency, 1M+ token context window. \\
Structured Output & Instructor + Pydantic & Enforces strict data models so frontend never crashes. \\
Frontend & Vanilla JS + Tailwind CSS & Zero build tools required (no heavy Node.js/Webpack). \\
Graph Rendering & Vis.js (Canvas) & Interactive, draggable knowledge graph visualization. \\
Slide Generation & \texttt{python-pptx} & Generates native Microsoft PowerPoint files on the fly. \\
Session Storage & SQLite (\texttt{sessions.db}) & Lightweight dual-write persistence for session history. \\
Vector Memory & HNSWlib + BAAI/bge-m3 & Fast semantic retrieval with forgetting curve decay. \\
\bottomrule
\end{tabular}
\end{table}

% -----------------------------------------------------------------------------
% Section 5: Top 10 Viva Questions & Winning Answers
% -----------------------------------------------------------------------------
\section{Top 10 Viva Questions with Winning Answers}

\subsection*{1. What are the ``Pedagogies'' or ``Skills'' in GAI-PLA?}
\begin{answerbox}
Teaching styles are not hardcoded in Python; they are pluggable Markdown files in \texttt{skills/}:
\begin{itemize}[noitemsep]
    \item \textbf{Feynman Technique:} Explains concepts in ultra-simple words using real-world analogies (no confusing jargon).
    \item \textbf{Socratic Method:} Never gives the answer directly; asks guiding questions so the student deduces the answer themselves.
    \item \textbf{Scaffolding:} Breaks complex math into small, supportive baby steps and gradually lets the student solve on their own.
\end{itemize}
\end{answerbox}

\subsection*{2. How does the frontend communicate with the backend?}
\begin{answerbox}
Through two standard web protocols:
\begin{enumerate}[noitemsep]
    \item \textbf{Standard REST API (\texttt{fetch}):} For submitting quiz answers, changing settings, and fetching session state (\texttt{GET /api/chat/\{sid\}/state}).
    \item \textbf{Server-Sent Events (SSE):} For streaming AI responses token-by-token (in chat and lessons) so the user doesn't wait 10 seconds staring at a blank screen.
\end{enumerate}
\end{answerbox}

\subsection*{3. How do you prevent AI hallucinations?}
\begin{answerbox}
We use a 3-layer guardrail:
\begin{enumerate}[noitemsep]
    \item \textbf{Jinja2 Prompt Templates:} All prompts live in structured \texttt{.j2} files with strict constraints.
    \item \textbf{Instructor Schema Enforcement:} The model cannot hallucinate arbitrary text when forced to output a rigid Pydantic schema.
    \item \textbf{Grounded Knowledge Nodes:} Questions are strictly scoped to verified nodes generated in Step 1.
\end{enumerate}
\end{answerbox}

\subsection*{4. What happens if the student refreshes the browser page?}
\begin{answerbox}
Nothing is lost! The backend has a dual-write store (\texttt{store.py}) that persists state to memory and an on-disk SQLite database (\texttt{sessions.db}). When the page reloads, the frontend calls \texttt{GET /api/chat/\{sid\}/state} and completely reconstructs the knowledge graph, diagnosis scores, and lessons.
\end{answerbox}

\subsection*{5. How does the PowerPoint (.pptx) export work?}
\begin{answerbox}
The AI first creates a slide outline, then layouts each slide in a fixed $1000 \times 562.5$ coordinate canvas (standard 16:9 widescreen). A 5-stage post-processing pipeline repairs any malformed JSON or LaTeX, and \texttt{src/export/pptx\_export.py} translates the elements directly into native PowerPoint shapes, text boxes, and tables using \texttt{python-pptx}.
\end{answerbox}

\subsection*{6. What is the ``Ebbinghaus Forgetting Curve'' in your project?}
\begin{answerbox}
German psychologist Hermann Ebbinghaus proved humans forget information exponentially over time unless reviewed. Our memory module implements this mathematical decay ($R = e^{-t/S}$). The system computes how long ago a student reviewed a weak point and schedules smart revision questions at the exact optimal interval.
\end{answerbox}

\subsection*{7. What was your individual contribution to the project?}
\begin{answerbox}
\textbf{Highlight these points confidently:}
\begin{itemize}[noitemsep]
    \item Adapted the codebase into the \textbf{GAI-PLA} personalized learning framework.
    \item Configured Google Gemini API via an OpenAI-compatible structured gateway.
    \item Integrated Windows execution and environment dependencies.
    \item Rebranded and customized the full user interface across diagnosis, chat, lessons, and slides.
    \item Conducted experimental validation on STEM topics (Pythagoras theorem case study).
\end{itemize}
\end{answerbox}

\subsection*{8. What are the limitations and future scope?}
\begin{answerbox}
\begin{itemize}[noitemsep]
    \item \textbf{Current Limitation:} Requires an active internet connection to call the Gemini API gateway.
    \item \textbf{Future Scope:} Running local quantized SLMs (Small Language Models) completely offline, adding voice conversation, and generating interactive 3D simulations.
\end{itemize}
\end{answerbox}

% -----------------------------------------------------------------------------
% Viva Strategy Tips
% -----------------------------------------------------------------------------
\vspace{0.4cm}
\begin{tipbox}
\textbf{How to handle the Viva presentation:}
\begin{enumerate}[noitemsep]
    \item \textbf{Lead with the Live Demo:} Show the Pythagoras theorem graph $\to$ 50\% diagnosis $\to$ generated slides. Seeing working software immediately impresses evaluators.
    \item \textbf{Speak about Architecture, not just Prompts:} Emphasize that GAI-PLA is a \textbf{State Machine System} (LangGraph + FastAPI), not just an API wrapper.
    \item \textbf{Stay calm on failures:} If an API rate limit or network lag occurs, explain: \textit{``This is the external cloud LLM gateway's latency; the local state machine and SQLite persistence remain intact.''}
\end{enumerate}
\end{tipbox}

\end{document}
